\documentclass[12pt]{article}
\usepackage[margin=1in]{geometry}
\usepackage{enumitem}
\usepackage{titlesec}
\usepackage{parskip}
\usepackage{graphicx}
\usepackage{xcolor}
\usepackage{hyperref}
\usepackage{fontenc}

\title{\textbf{Posner 2004 RG Notes}\\
\large Posner, Daniel N.. ``The Political Salience of Cultural Difference: Why Chewas and Tumbukas Are Allies in Zambia and Adversaries in Malawi'' [2004].\\
\textit{American Political Science Review} 98(4): 529-545.}
\author{}
\date{}

\begin{document}

\maketitle

When do cultural cleavages become politically salient? See Lipset \& Rokkan (1967) who introduce three such conditions: empirical framework, normative element, and organizational framework

Claims that the salience of cultural cleavages can be largely determined by contextual factors \textit{external to} the cleavages themselves: namely, whether or not the resultant groups formed from those cleavages are effective vehicles for political mobilization

Single biggest factor is size: are the resulting groups large enough such that they can be meaningful political groups?

Natural experiment with Chewa and Tumbuka tribes in Zambia/Malawi: cleavage becomes politically salient/antagonistic in Malawi, but not so in Zambia -- much more of an alliance. Conducted surveys within villages straddling the Malawi/Zambian border to evaluate the differential political impacts of those cleavages across the border.

Potentially interesting note here -- using colonial borders as a framework for a natural experiment, since they were drawn administratively rather than culturally and are therefore largely random

\end{document}
